\begin{center}
\textbf{\large{Tóm tắt}	}
\end{center}

\addcontentsline{toc}{chapter}{Tóm tắt}

\begin{small}
\textbf{\textit{Tóm tắt:}} Thực khuẩn thể là vi-rút xâm nhiễm vi khuẩn, được xem là một giải pháp tiềm năng thay thế kháng sinh thông qua liệu pháp thực khuẩn thể. Trong ứng dụng trị liệu, việc xác định chính xác lối sống của thực khuẩn thể, độc lực hay ôn hòa, là một yêu cầu quan trọng, bởi phần lớn các ứng dụng điều trị ưu tiên thực khuẩn thể độc lực; thực khuẩn thể ôn hòa có thể tích hợp vật liệu di truyền vào hệ gen vật chủ theo chu trình tiềm tan và góp phần lan truyền các gen độc lực hay gen kháng kháng sinh qua chuyển gen ngang, từ đó làm tăng rủi ro cho điều trị. Các phương pháp truyền thống dựa trên nuôi cấy trong phòng thí nghiệm thường tốn nhiều thời gian và nguồn lực, khó áp dụng cho phần lớn thực khuẩn thể chưa được nuôi cấy hay gán nhãn, đặt ra nhu cầu về các phương pháp phân loại dựa trên tính toán. Các công cụ hiện có như DeePhage, PhaTYP và ProkBERT đạt hiệu suất khả quan nhưng vẫn mang trong mình những nhược điểm riêng.

Luận văn này đề xuất PhaBERT, một kiến trúc học sâu kết hợp mô hình nền tảng hệ gen DNABERT-2 với mô-đun tích chập đa nhân MKC. Nhờ đó, PhaBERT tận dụng tri thức tiền huấn luyện trên dữ liệu di truyền đa loài cùng với khả năng hoạt động trực tiếp trên dữ liệu di truyền thô thông qua mã hóa BPE mà không bị rò rỉ thông tin.

Thực nghiệm trên bốn nhóm độ dài contig (100-400 bp, 400-800 bp, 800-1.200 bp, 1.200-1.800 bp) với kiểm định chéo phân tầng 5 phần cho thấy PhaBERT đạt độ chính xác lần lượt 82,01\%, 87,33\%, 89,90\% và 91,38\%, cải thiện 0,36--6,65\% so với các phương pháp tiên tiến gồm PhaTYP, DeePhage và ProkBERT. Nghiên cứu loại bỏ thành phần cho thấy mô-đun MKC đóng góp cải thiện độ chính xác từ 0,49\% đến 3,02\% so với đường cơ sở DNABERT-2 chỉ dùng đầu phân loại tuyến tính, với mức cải thiện tăng dần theo độ dài contig. Mô hình duy trì sự cân bằng giữa độ nhạy (82,07\%--92,38\%) và độ đặc hiệu (81,86\%--87,63\%) trên các phạm vi độ dài được khảo sát. Kết quả nghiên cứu cho thấy hiệu quả của việc kết hợp mô hình nền tảng hệ gen với kiến trúc đặc thù theo tác vụ. Đây là một hướng tiếp cận có nhiều triển vọng cho các bài toán phân tích chuỗi sinh học, nhất là trong phân tích dữ liệu hệ gen môi trường với các chuỗi không hoàn chỉnh hoặc bị phân mảnh.

\vspace*{1cm}
\textbf{\textit{Từ khóa}}: \textit{Thực khuẩn thể}, \textit{Phân loại lối sống}, \textit{Học sâu}, \textit{BERT}, \textit{Mạng nơ-ron tích chập}, \textit{Tin sinh học}, \textit{Hệ gen môi trường}

\end{small}
